\section{研究总结}

本文围绕轻量化图像识别模型设计这一主题，完成了从问题分析、相关技术梳理、模型方案设计到实验验证的完整研究流程。主要工作包括：分析图像识别任务的研究背景与应用需求；总结卷积神经网络、残差学习和注意力机制等关键技术；在基线模型基础上设计融合通道重标定和多尺度特征聚合的改进方法；通过对比实验与消融实验验证所提方法的有效性。

实验结果表明，本文方法能够在较小额外开销下提升分类准确率，说明轻量化结构优化是一条可行的改进路径。对于本科毕业论文写作而言，上述研究过程也展示了规范论文应具备的完整逻辑链条，即“提出问题、分析问题、设计方法、验证效果、总结结论”。

\section{不足与展望}

尽管本文取得了一定结果，但仍存在以下不足：

\begin{enumerate}
  \item 实验数据集规模有限，对更复杂真实场景的覆盖仍然不足。
  \item 模型改进主要集中在结构层面，对损失函数设计和样本不平衡处理考虑较少。
  \item 本文尚未在移动端或嵌入式设备上完成完整部署验证，缺乏端侧性能测试。
\end{enumerate}

未来可以从以下几个方向继续开展研究：

\begin{enumerate}
  \item 引入更大规模或更贴近实际应用的数据集，进一步验证模型泛化能力。
  \item 结合知识蒸馏、剪枝和量化等技术，继续压缩模型体积并提升推理效率。
  \item 将当前方法扩展到目标检测、缺陷识别或医学影像分析等更具应用价值的视觉任务中。
\end{enumerate}

\section{结语}

本科毕业设计（论文）不仅要求完成一定的研究工作，更要求在表达层面符合学术规范。一个稳定、易维护且贴近学校格式要求的 LaTeX 模板，能够帮助学生把更多精力投入到研究内容本身。希望本模板既能作为论文写作工具，也能成为后续同学开展学术表达训练的起点。
